% ID: SEL-D0066-OA-E03
% ESTADO: cotejada. Segunda pasada visual y borrador compilado revisados.
% Fuente: Selectividad/sel_2010_matematicas_aplicadas/mat_apli/mat_apli_sept/mat_apli_sept_rb/EXAMEN-ANDALUCÍA-2.pdf
% SHA-256: 24c494a5d013dbc7f5adaac9896dee2d79f22540518359c67eac29945a26d4f0
% Página: 1; opción A; ejercicio original 3.

Una empresa utiliza dos servidores para conectarse a Internet. El primero, $S_1$, lo utiliza el $45\%$ de las veces y el segundo, $S_2$, el resto.

Cuando se conecta a Internet con $S_1$, los ordenadores se bloquean el $5\%$ de las veces, y cuando lo hace con $S_2$ el $8\%$. Si un día, al azar, la empresa está conectada a Internet,
\begin{enumerate}[label=\alph*)]
\item \textbf{(1.25 puntos)} \textquestiondown{}cuál es la probabilidad de que los ordenadores se queden bloqueados?
\item \textbf{(1.25 puntos)} \textquestiondown{}Cuál es la probabilidad de que la empresa esté utilizando el servidor $S_1$, sabiendo que los ordenadores se han quedado bloqueados?
\end{enumerate}
